The connection is quite natural, and I think it gives a useful geometric organization of the manifold version of Nelson mechanics with electromagnetic fields.

The central idea is:

$$ \boxed{ \text{Riemannian geometry} + \text{stochastic parallel transport} + U(1)\text{ connection} } $$

rather than trying to put everything into a single "velocity potential."

Nelson's original framework already admits a Lagrangian with both scalar and vector potentials,

$$ L=\frac12 m_{ij}v^iv^j-\phi+A_iv^i, $$

as Nelson himself emphasized in his later discussion of stochastic mechanics.

1. Start with Brownian motion on a manifold

Let \((M,g)\) be the configuration manifold.

The natural diffusion is not

$$ dq^i=b^i\,dt+\sqrt{\frac{\hbar}{m}}\,dW^i $$

written in arbitrary coordinates, because \(dW^i\) does not transform as a tangent vector.

Instead one uses the Levi-Civita connection and stochastic parallel transport.

A convenient Stratonovich form is

$$ \boxed{ dX_t=b(X_t,t)\,dt + \sqrt{\frac{\hbar}{m}}\, E_a(X_t)\circ dW_t^a , } $$

where

$$ g(E_a,E_b)=\delta_{ab}. $$

The resulting generator contains the Laplace--Beltrami operator,

$$ \mathcal L = b^i\nabla_i + \frac{\hbar}{2m}\Delta_g. $$

This is the geometric replacement for ordinary Euclidean Brownian motion.

The Eells--Elworthy--Malliavin construction makes precisely this use of stochastic parallel transport; it constructs Brownian motion intrinsically on a Riemannian manifold.

2. Nelson's two velocities

For a diffusion \(X_t\), Nelson introduces forward and backward derivatives

$$ D X_t, \qquad D_*X_t. $$

On a manifold it is better to regard these as tangent vectors:

$$ b^i_+=DX^i, \qquad b^i_-=D_*X^i. $$

Define

$$ v = \frac12(b_++b_-), $$

and

$$ u = \frac12(b_+-b_-). $$

They are respectively the current velocity and osmotic velocity.

For a density \(\rho\),

$$ \boxed{ u = \frac{\hbar}{2m}\nabla\ln\rho } $$

in the standard Nelson normalization.

The continuity equation is

$$ \boxed{ \partial_t\rho+\operatorname{div}(\rho v)=0. } $$

The important point is that this construction is intrinsically meaningful on \(M\).

3. Now introduce electromagnetism

Introduce a \(U(1)\) connection

$$ A=A_i\,dq^i. $$

Its curvature is

$$ \boxed{ F=dA, \qquad F_{ij} = \partial_iA_j-\partial_jA_i. } $$

The classical Lagrangian is

$$ L = \frac{m}{2}g_{ij}\dot q^i\dot q^j + qA_i\dot q^i - q\phi. $$

The corresponding equation is

$$ m\nabla_{\dot q}\dot q = q\,\iota_{\dot q}F -q\,d\phi. $$

In three-dimensional language,

$$ m\dot v = q(E+v\times B). $$

The key geometric observation is:

$$ \boxed{ A\text{ is a connection},\qquad F\text{ is its curvature}. } $$

Therefore the magnetic force is already a curvature effect.

4. Put the connection into Nelson mechanics

There are two closely related ways to do this.

Route A: modify the stochastic acceleration

Keep the Brownian motion on \(M\), but replace the Nelson dynamical equation by

$$ \boxed{ \frac{m}{2} \left(D D_*+D_*D\right)X = -q\,d\phi + q\,\iota_vF. } $$

The left side is the stochastic analogue of acceleration.

The right side contains the electromagnetic force.

This is the most direct Nelson formulation.

Route B: introduce a line bundle

This is geometrically more powerful.

Let

$$ L\rightarrow M $$

be a complex line bundle with \(U(1)\) connection

$$ \nabla^A = d-\frac{iq}{\hbar}A. $$

Then the covariant derivative is

$$ \boxed{ \nabla_i^A = \nabla_i-\frac{iq}{\hbar}A_i. } $$

The wavefunction is no longer merely a complex scalar globally; it is naturally a section

$$ \psi\in\Gamma(L). $$

The magnetic Schrödinger operator becomes

$$ \boxed{ H = \frac{1}{2m} (-i\hbar\nabla^A)^* (-i\hbar\nabla^A) + q\phi. } $$

Locally,

$$ H = \frac{1}{2m} \left(-i\hbar\nabla-qA\right)^2 + q\phi. $$

This is exactly the bundle structure that becomes important when \(A\) cannot be represented globally by a single one-form.

5. The phase becomes a connection statement

Write locally

$$ \psi=\sqrt{\rho}\,e^{iS/\hbar}. $$

Then the physically relevant current velocity is not simply

$$ \frac{1}{m}\nabla S. $$

It is

$$ \boxed{ v = \frac{1}{m} \left( dS-qA \right)^\sharp . } $$

This is crucial.

Taking the exterior derivative,

$$ dv^\flat = -\frac{q}{m}F $$

locally, because \(d^2S=0\).

Therefore

$$ \boxed{ \nabla\times v = -\frac qm B } $$

in Euclidean three-space, with the appropriate sign depending on the convention for \(A\).

So magnetic vorticity is not an obstruction to the geometric formulation.

It is precisely what one expects from a nontrivial connection.

This is very relevant to your earlier concern that Nelson/Hamilton--Jacobi formulations can fail when the velocity is not a gradient.

Without electromagnetism,

$$ v^\flat=\frac1m dS $$

would imply

$$ dv^\flat=0. $$

With a gauge field,

$$ v^\flat=\frac1m(dS-qA), $$

and therefore

$$ dv^\flat=-\frac qmF. $$

Thus the correct generalization is not

$$ v=\nabla S/m, $$

but

$$ \boxed{ mv^\flat+qA=dS. } $$
6. This solves an important conceptual problem

There are two different sources of non-gradient behavior:

Topological/quantum circulation

Even with \(F=0\),

$$ \oint v^\flat\neq0 $$

can occur when the configuration space is not simply connected.

This is the familiar circulation/holonomy issue.

Electromagnetic curvature

When

$$ F\neq0, $$

even locally,

$$ dv^\flat\neq0. $$

Thus

$$ \boxed{ \text{non-gradient velocity} = \text{topology and/or gauge curvature}. } $$

This is much more satisfactory than treating vorticity merely as a pathology of the Madelung representation.

7. Stochastic parallel transport enters naturally

Now suppose \(X_t\) is the stochastic trajectory.

There are actually two parallel transports:

Tangent-bundle transport

The Levi-Civita connection transports tangent vectors,

$$ P^{LC}_{0,t}:T_{X_0}M\rightarrow T_{X_t}M. $$

This is needed to compare stochastic velocities at different points.

Gauge transport

The \(U(1)\) connection transports complex phases,

$$ U^A_{0,t}:L_{X_0}\rightarrow L_{X_t}. $$

Locally,

$$ \boxed{ U^A[X] = \exp\left( \frac{iq}{\hbar} \int_X A \right) } $$

with the appropriate path-ordering trivial for \(U(1)\).

Around a closed loop,

$$ U^A[\gamma] = \exp\left( \frac{iq}{\hbar} \oint_\gamma A \right). $$

By Stokes' theorem,

$$ \oint_\gamma A = \int_\Sigma F. $$

Therefore

$$ \boxed{ \text{holonomy} = \exp\left( \frac{iq}{\hbar} \int_\Sigma F \right). } $$

This is the stochastic version of the Aharonov--Bohm mechanism.

8. Now the Kaluza picture becomes useful

The Kaluza metric can be written schematically as

$$ ds^2_{K} = g_{ij}dq^i dq^j + R^2 \left( dy+\frac{q}{\hbar}A_i dq^i \right)^2. $$

The one-form

$$ \boxed{ \theta=dy+\frac{q}{\hbar}A } $$

is the connection form of the \(U(1)\) bundle.

Its curvature is

$$ d\theta = \frac{q}{\hbar}F. $$

Thus the Kaluza lift says:

$$ \boxed{ A\rightarrow\text{geometry of the fiber bundle}. } $$

Now stochastic mechanics can be performed on the total space.

The stochastic process has a horizontal component on \(M\) and a vertical component along the \(U(1)\) fiber.

Horizontal stochastic transport is defined by

$$ \theta(dX)=0. $$

Hence

$$ dy = -\frac{q}{\hbar}A_i\,dq^i. $$

Integrating,

$$ y_t-y_0 = -\frac{q}{\hbar}\int A. $$

So the vertical Kaluza coordinate records the accumulated gauge phase.

That is a remarkably clean correspondence:

$$ \boxed{ \text{Kaluza fiber coordinate} \sim \text{gauge phase}. } $$
9. The complete geometric structure

We can therefore put the whole construction together:

$$ \boxed{ \begin{array}{ccc} (M,g) & & U(1)\text{ bundle }L\to M \\ \downarrow && \downarrow \\ \text{Brownian motion} && \text{gauge parallel transport} \\ \downarrow && \downarrow \\ \text{Nelson }v,u && A,F,\text{ holonomy} \\ \downarrow && \downarrow \\ \text{stochastic acceleration} && \text{Lorentz force} \\ \multicolumn{3}{c}{ \Downarrow }\\ \multicolumn{3}{c}{ \text{magnetic Schrödinger equation} } \end{array} } $$

The roles are nicely separated:

$$ \boxed{ \begin{aligned} g &\quad\text{determines the Brownian geometry},\\ A &\quad\text{determines phase transport},\\ F=dA &\quad\text{determines magnetic curvature},\\ \rho &\quad\text{determines osmotic velocity},\\ S &\quad\text{determines the gauge-covariant current}. \end{aligned}} $$
10. Where stochastic parallel transport is more than a mathematical convenience

This is where I think the construction becomes particularly relevant to your proposed "true Nelson with spin" program.

For a scalar particle, the \(U(1)\) bundle is enough.

For spin, however, the relevant bundle is no longer a line bundle. One needs a spinor bundle

$$ S(M)\rightarrow M $$

with a connection

$$ \nabla^{S,A} = \nabla^{LC} + \frac{iq}{\hbar}A + \text{spin connection}. $$

Stochastic parallel transport then acts on spinors:

$$ \boxed{ U_t: S_{X_0}M\rightarrow S_{X_t}M. } $$

The Levi-Civita part transports the local frame; the \(U(1)\) part transports electromagnetic phase; and the spin connection transports the spin degree of freedom.

Schematically,

$$ \boxed{ \text{stochastic trajectory} + \text{frame transport} + \text{gauge transport} \rightarrow \text{spinor-valued stochastic process}. } $$

This is substantially more promising for a genuine stochastic derivation of the Pauli equation than simply inserting a known Pauli wavefunction into an SDE.

11. Curvature should then generate the spin coupling

For a spinor connection,

$$ [\nabla_\mu,\nabla_\nu]\psi = \frac14R_{\mu\nu ab}\gamma^{ab}\psi + \frac{iq}{\hbar}F_{\mu\nu}\psi. $$

So there are two curvatures:

$$ \boxed{ R_{\mu\nu ab} \quad\text{and}\quad F_{\mu\nu}. } $$

When one squares the Dirac operator, the electromagnetic curvature produces the Pauli term

$$ \sim \frac{q\hbar}{2m}\,\boldsymbol{\sigma}\cdot\mathbf B. $$

This is precisely the geometric mechanism you were interested in earlier:

$$ \boxed{ \text{parallel transport} \rightarrow \text{curvature} \rightarrow \sigma\cdot B. } $$

The important conceptual point is that spin does not have to be introduced as an ad hoc internal classical variable. It can arise from stochastic transport in a spin bundle.

12. And this connects back to Franchi--Le Jan

For ordinary nonrelativistic Nelson mechanics the base space is Riemannian, and Brownian motion can be constructed using the Levi-Civita connection.

For a relativistic extension, ordinary Brownian motion on spacetime is problematic because spacetime is Lorentzian. Franchi and Le Jan instead construct relativistic diffusions on the unit future tangent bundle, with the velocity undergoing hyperbolic Brownian motion and being transported along the spacetime trajectory.

That is extremely relevant to your proposed relativistic Nelson program.

The natural hierarchy becomes

$$ \boxed{ \begin{array}{c} \text{configuration manifold }M\\ \downarrow\\ TM\text{ / frame bundle}\\ \downarrow\\ \text{stochastic parallel transport}\\ \downarrow\\ \text{gauge bundle}\\ \downarrow\\ \text{spin bundle}\\ \downarrow\\ \text{relativistic stochastic dynamics}. \end{array}} $$

Franchi--Le Jan provides an established geometric model for the relativistic stochastic part; their construction explicitly uses parallel transport of velocities along the random spacetime curves.

13. The synthesis I would pursue

For your SHP–Nelson direction, I would not start by trying to write a gigantic "relativistic Nelson equation."

I would build it in layers:

$$ \boxed{ \begin{array}{rcl} (M,g) &\longrightarrow& \text{Brownian diffusion} \\ &&\downarrow\\ (M,g,A) &\longrightarrow& \text{gauge-covariant Nelson mechanics} \\ &&\downarrow\\ (P_{\mathrm{Spin}},\omega,A) &\longrightarrow& \text{spin stochastic transport} \\ &&\downarrow\\ \text{Lorentzian }(M,g) &\longrightarrow& \text{Franchi--Le Jan type relativistic diffusion} \\ &&\downarrow\\ \text{SHP extended spacetime} &\longrightarrow& \text{relativistic stochastic dynamics}. \end{array}} $$

The key structural lesson is:

$$ \boxed{ \textbf{Do not identify the stochastic velocity with a gradient.} } $$

Instead,

$$ \boxed{ m v^\flat=dS-qA } $$

is the gauge-covariant relation, while

$$ \boxed{ F=dA } $$

measures its local failure to be irrotational.

Then stochastic parallel transport gives the missing global structure, and for spin the transport is promoted from the tangent bundle/\(U(1)\) line bundle to the spin bundle.

That, in my view, is the clean geometric route from Nelson → magnetic Nelson → bundle-valued Nelson → spin → relativistic stochastic mechanics, rather than the less satisfactory procedure of starting with a Schrödinger/Pauli/Dirac solution and constructing an SDE from it afterward. Nelson himself later emphasized that gauge theories may be essential to a relativistic stochastic program.
